\documentclass[12pt,a4paper]{article}
\usepackage[utf8]{inputenc}
\usepackage[T2A]{fontenc}
\usepackage[russian]{babel}
\usepackage{amsmath, amsfonts, amssymb, amsthm}
\usepackage{geometry}
\usepackage{graphicx}
\usepackage{enumitem}

\geometry{left=2.5cm, right=1.5cm, top=2cm, bottom=2cm}

\newtheorem{theorem}{Теорема}
\newtheorem{definition}{Определение}
\renewcommand{\qedsymbol}{$\blacksquare$}

\title{Билет №92 \\ \small Определение многомерных БИХ систем. Основные подходы к синтезу БИХ систем. Параллелизация алгоритмов вычисления в БИХ системах. (СпБПУ)}
\author{Сериков Владислав Александрович}
\date{16 мая 2026}

\begin{document}

\maketitle
\tableofcontents
\newpage

\section{Определение многомерных БИХ-систем}

\begin{definition}
\textbf{Дискретная система с бесконечной импульсной характеристикой (БИХ-система)} — это класс линейных стационарных дискретных систем, реакция которых на единичный импульс имеет бесконечную длительность. В отличие от КИХ-систем, БИХ-системы содержат обратную связь (рекурсивную часть).
\end{definition}

\begin{definition}
\textbf{Многомерная БИХ-система ($m$-мерная)} — это система обработки многомерных дискретных сигналов $x[n_1, n_2, \dots, n_m]$, которая описывается многомерным линейным разностным уравнением с постоянными коэффициентами, включающим рекурсивные члены.
\end{definition}

Для простоты и наглядности наиболее часто рассматривают двумерные (2D) системы ($m=2$), которые широко применяются в обработке изображений и пространственной фильтрации. Разностное уравнение каузальной двумерной БИХ-системы имеет вид:
\begin{equation}
    y[n_1, n_2] = \sum_{k_1=0}^{M_1} \sum_{k_2=0}^{M_2} a_{k_1, k_2} x[n_1-k_1, n_2-k_2] - \sum_{\substack{l_1=0 \\ (l_1, l_2) \neq (0,0)}}^{N_1} \sum_{l_2=0}^{N_2} b_{l_1, l_2} y[n_1-l_1, n_2-l_2],
\end{equation}
где $x[n_1, n_2]$ — входной сигнал, $y[n_1, n_2]$ — выходной сигнал, $a_{k_1, k_2}$ и $b_{l_1, l_2}$ — коэффициенты прямой и обратной связей соответственно. Предполагается, что $b_{0,0} = 1$.

Применяя двумерное Z-преобразование, получаем передаточную функцию системы:
\begin{equation}
    H(z_1, z_2) = \frac{Y(z_1, z_2)}{X(z_1, z_2)} = \frac{A(z_1, z_2)}{B(z_1, z_2)} = \frac{\sum_{k_1=0}^{M_1} \sum_{k_2=0}^{M_2} a_{k_1, k_2} z_1^{-k_1} z_2^{-k_2}}{\sum_{l_1=0}^{N_1} \sum_{l_2=0}^{N_2} b_{l_1, l_2} z_1^{-l_1} z_2^{-l_2}}.
\end{equation}

Ключевой проблемой многомерных БИХ-систем является обеспечение их устойчивости (BIBO — Bounded Input Bounded Output), так как многомерные полиномы, в отличие от одномерных, в общем случае не факторизуются на множители первого порядка.

\begin{theorem}[Теорема Хуанга об устойчивости 2D БИХ-систем]
Каузальная двумерная рекурсивная система с передаточной функцией $H(z_1, z_2) = \frac{A(z_1, z_2)}{B(z_1, z_2)}$, где $A$ и $B$ — взаимно простые полиномы от $z_1^{-1}$ и $z_2^{-1}$, является устойчивой (BIBO) тогда и только тогда, когда выполняются два условия (для положительных степеней комплексных переменных $z_i$):
\begin{enumerate}
    \item $B(z_1, \infty) \neq 0$ для всех $|z_1| \ge 1$;
    \item $B(z_1, z_2) \neq 0$ для всех $|z_1| = 1$ и $|z_2| \ge 1$.
\end{enumerate}
\end{theorem}

\begin{proof}
Доказательство базируется на принципе непрерывности корней многочлена.
Введем обозначения в положительных степенях. Устойчивость требует отсутствия нулей знаменателя в области $|z_1| \ge 1, |z_2| \ge 1$.

\textit{Необходимость.} Если система устойчива, то $B(z_1, z_2) \neq 0$ в области $|z_1| \ge 1 \cap |z_2| \ge 1$. Очевидно, что подмножество этой области при $|z_1| = 1$ удовлетворяет условию 2. Устремляя $z_2 \to \infty$ при фиксированном $|z_1| \ge 1$, получаем выполнение условия 1.

\textit{Достаточность.} Предположим, условия 1 и 2 выполнены. Рассмотрим $B(z_1, z_2) = 0$ как алгебраическое уравнение относительно $z_2$ при заданном параметре $z_1$. Согласно условию 2, если $|z_1|=1$, то все корни $z_2$ лежат внутри единичного круга ($|z_2| < 1$).
При непрерывном изменении $z_1$ в области $|z_1| \ge 1$, корни $z_2(z_1)$ также изменяются непрерывно.
Условие 1 ($B(z_1, \infty) \neq 0$) гарантирует, что степень полинома по $z_2$ не уменьшается при любом $|z_1| \ge 1$, следовательно, корни не могут «приходить» из бесконечности.
Для того чтобы при каком-либо $|z_1^*| \ge 1$ корень $z_2$ оказался в области $|z_2| \ge 1$, он должен пересечь границу $|z_2| = 1$ в силу непрерывности. Но согласно условию 2, на границе $|z_1|=1$ пересечение невозможно, а аналитическое продолжение гарантирует, что корни остаются в области $|z_2| < 1$ для всех $|z_1| \ge 1$. Следовательно, $B(z_1, z_2) \neq 0$ для $|z_1| \ge 1, |z_2| \ge 1$.
\end{proof}


\section{Основные подходы к синтезу БИХ-систем}

Синтез БИХ-системы заключается в нахождении коэффициентов $a$ и $b$ передаточной функции $H(z)$, удовлетворяющей заданным частотным или временным характеристикам.

\subsection{Одномерный синтез: Метод билинейного преобразования}
Является основным методом проектирования одномерных БИХ-фильтров на основе аналоговых прототипов (Баттерворта, Чебышева, эллиптических).

\textbf{Алгоритм синтеза методом билинейного преобразования:}
\begin{enumerate}[label=\textbf{Шаг \arabic*.}, leftmargin=*]
    \item \textbf{Спецификация требований:} Задаются требования к цифровому фильтру: частота пропускания $\omega_p$, частота задерживания $\omega_s$, неравномерность в полосе пропускания $R_p$ и минимальное затухание $A_s$.
    \item \textbf{Искажение частотной шкалы (Pre-warping):} Пересчет цифровых частот в частоты аналогового прототипа по формуле: $\Omega = \frac{2}{T} \tan\left(\frac{\omega}{2}\right)$, где $T$ — период дискретизации (часто принимают $T=2$).
    \item \textbf{Синтез аналогового прототипа:} Расчет передаточной функции аналогового фильтра $H_a(s)$, удовлетворяющего требованиям $\Omega_p, \Omega_s, R_p, A_s$.
    \item \textbf{Билинейная подстановка:} Замена переменной Лапласа $s$ на дискретную переменную $z$ по формуле $s = \frac{2}{T}\frac{1 - z^{-1}}{1 + z^{-1}}$. Получение $H(z) = \left. H_a(s) \right|_{s = \frac{1 - z^{-1}}{1 + z^{-1}}}$.
\end{enumerate}

\textbf{Минимальный пример:}
Требуется синтезировать ФНЧ 1-го порядка.
Аналоговый прототип Баттерворта 1-го порядка с частотой среза $\Omega_c=1$: $H_a(s) = \frac{1}{s+1}$.
Применяем подстановку (с $T=2$):
$H(z) = \frac{1}{\frac{1-z^{-1}}{1+z^{-1}} + 1} = \frac{1+z^{-1}}{2}$. (Получен простейший ФНЧ).

\subsection{Многомерный синтез}
Из-за невозможности факторизации многочленов, синтез многомерных БИХ-систем существенно сложнее.

\begin{itemize}
    \item \textbf{Метод преобразования 1D в 2D:} Использование преобразований (например, обобщенного преобразования МакКлеллана) для перевода стабильных одномерных БИХ-фильтров в многомерные.
    \item \textbf{Метод сепарабельных фильтров:} Проектирование фильтра как каскада двух независимых 1D фильтров: $H(z_1, z_2) = H_1(z_1)H_2(z_2)$. Значительно упрощает синтез и проверку устойчивости, но ограничивает форму частотного отклика (только прямоугольные границы).
    \item \textbf{Оптимизационные методы (Пространственно-частотный синтез):} Формируется целевая функция ошибки между желаемой АЧХ $H_d(\omega_1, \omega_2)$ и реальной $H(\omega_1, \omega_2)$. Процесс сводится к нелинейной оптимизации коэффициентов $a_{k_1, k_2}$ и $b_{l_1, l_2}$ с ограничениями на устойчивость. Часто применяют методы генетических алгоритмов или метод наименьших квадратов.
\end{itemize}


\section{Параллелизация алгоритмов вычисления в БИХ-системах}

Параллельная обработка в КИХ-системах достигается тривиально (сигнал можно разделить на независимые блоки). В БИХ-системах обратная связь $y[n] = f(y[n-1])$ создает жесткую \textit{временную зависимость (iteration bound)}, что препятствует прямой конвейеризации и распараллеливанию.

\subsection{Метод опережающих вычислений (Look-Ahead)}

Метод трансформирует структуру системы таким образом, чтобы текущий выход зависел не от предыдущего $y[n-1]$, а от $y[n-M]$, где $M$ — уровень параллелизма. Это освобождает $M$ тактов для введения конвейерных регистров.

\begin{theorem}[Теорема эквивалентности при опережающих вычислениях]
Линейное рекурсивное уравнение первого порядка $y[n] = a y[n-1] + x[n]$ математически эквивалентно уравнению
\begin{equation}
    y[n] = a^M y[n-M] + \sum_{i=0}^{M-1} a^i x[n-i]
\end{equation}
для любого целого $M \ge 1$.
\end{theorem}

\begin{proof}
Доказательство проведем методом математической индукции по $M$.
\newline
\textit{База индукции ($M=1$):} Уравнение принимает вид $y[n] = a^1 y[n-1] + a^0 x[n]$, что совпадает с исходным уравнением. База верна.
\newline
\textit{Шаг индукции:} Предположим, что утверждение верно для $M = k$:
$y[n] = a^k y[n-k] + \sum_{i=0}^{k-1} a^i x[n-i]$.
Докажем для $M = k+1$. Рассмотрим член $y[n-k]$ и выразим его из исходного уравнения, сдвинутого на $k$ тактов: $y[n-k] = a y[n-k-1] + x[n-k]$. Подставим это в предположение индукции:
\begin{align*}
    y[n] &= a^k \big( a y[n-k-1] + x[n-k] \big) + \sum_{i=0}^{k-1} a^i x[n-i] \\
         &= a^{k+1} y[n-(k+1)] + a^k x[n-k] + \sum_{i=0}^{k-1} a^i x[n-i] \\
         &= a^{k+1} y[n-(k+1)] + \sum_{i=0}^{k} a^i x[n-i].
\end{align*}
Полученное выражение совпадает с формулой для $M = k+1$. Теорема доказана.
\end{proof}

\textbf{Алгоритм конвейеризации (Scattered Look-Ahead):}
\begin{enumerate}[label=\textbf{Шаг \arabic*.}, leftmargin=*]
    \item Определить требуемый уровень параллелизма/конвейеризации $M$.
    \item Представить БИХ-систему в Z-области: $H(z) = \frac{1}{1 - a z^{-1}}$.
    \item Умножить числитель и знаменатель на полином $\sum_{i=0}^{M-1} a^i z^{-i}$.
    \item Получить эквивалентную систему $H(z) = \frac{\sum_{i=0}^{M-1} a^i z^{-i}}{1 - a^M z^{-M}}$.
    \item Реализовать вычисления: знаменатель обеспечивает рекурсию с задержкой на $M$ тактов, а числитель представляет собой КИХ-часть, которая легко распараллеливается.
\end{enumerate}

\textbf{Пример:} Для $M=2$ фильтр $H(z) = \frac{1}{1 - 0.5z^{-1}}$ преобразуется в $H(z) = \frac{1 + 0.5z^{-1}}{1 - 0.25z^{-2}}$. Теперь рекурсивный контур имеет задержку $z^{-2}$, что позволяет обрабатывать два отсчета параллельно.

\subsection{Блочная обработка (Block Processing)}

Блочная обработка вычисляет сразу блок из $L$ выходных отсчетов на основе блока из $L$ входных отсчетов, формулируя задачу в пространстве состояний.

\textbf{Алгоритм перевода в блочную структуру:}
\begin{enumerate}[label=\textbf{Шаг \arabic*.}, leftmargin=*]
    \item Описать систему в пространстве состояний:
    $x_{n+1} = A x_n + B u_n$, \quad $y_n = C x_n + D u_n$.
    \item Задать размер блока $L$.
    \item Рекурсивно подставить уравнения состояния для получения выражения для $x_{n+L}$:
    $x_{n+L} = A^L x_n + A^{L-1}B u_n + \dots + B u_{n+L-1}$.
    \item Сформировать векторные блочные входы $\mathbf{U}_k$ и выходы $\mathbf{Y}_k$.
    \item Реализовать матричные вычисления $\mathbf{X}_{k+1} = \mathcal{A} \mathbf{X}_k + \mathcal{B} \mathbf{U}_k$, где матричные операции умножения выполняются строго параллельно.
\end{enumerate}

Для \textbf{многомерных систем} параллелизация достигается за счет \textit{волновой (wave-front) обработки}. Поскольку $y[n_1, n_2]$ зависит от отсчетов слева и сверху, можно вычислять элементы, лежащие на побочной диагонали (антидиагонали) матрицы изображения одновременно, так как они не зависят друг от друга, формируя вычислительный волновой фронт, движущийся из верхнего левого в правый нижний угол.

\end{document}
